\section{Introducción}
\label{sec:introduccion}

Las plagas y enfermedades de los cultivos son una de las principales amenazas para la
seguridad alimentaria. Una evaluación global sobre cinco cultivos básicos estimó que las
pérdidas de rendimiento atribuibles a patógenos y plagas alcanzan, en el caso de la papa,
un 17,2\,\% en promedio (8,1--21,0\,\% según la región), y que las pérdidas más altas se
concentran en regiones con déficit de alimentos~\cite{savary2019burden}. El manejo de estas
enfermedades depende de un diagnóstico temprano, que todavía se apoya mayormente en la
inspección visual realizada por especialistas: un recurso escaso y costoso, sobre todo para
productores pequeños y alejados de los centros técnicos.

El aprendizaje profundo y la masificación de los teléfonos inteligentes abrieron la
posibilidad de automatizar ese diagnóstico a partir de una fotografía de la hoja. El dataset
PlantVillage~\cite{hughes2015plantvillage} fue el punto de partida de esta línea: sobre
54\,306 imágenes de 14 especies, una red convolucional alcanzó un 99,35\,\% de accuracy en
un conjunto de test separado~\cite{mohanty2016using}. Sin embargo, el mismo trabajo reportó
que el accuracy caía al 31,4\,\% sobre imágenes obtenidas de otras fuentes. Estudios
posteriores mostraron que PlantVillage contiene sesgos de captura que un modelo puede
explotar (por ejemplo, un clasificador entrenado con solo 8 píxeles del fondo alcanza un
49\,\% de accuracy~\cite{noyan2022uncovering}) y propusieron conjuntos de imágenes tomadas
en condiciones reales, como PlantDoc~\cite{singh2020plantdoc}. En consecuencia, un
accuracy cercano al 100\,\% en PlantVillage no garantiza por sí solo un modelo útil en el
campo, y su evaluación debe complementarse con controles de fuga de datos, explicabilidad y
pruebas fuera de la distribución de entrenamiento.

% --- Caso real de aplicación (<= 7 años). Propuesta: PlantVillage Nuru. ---------------
Un caso real que ilustra tanto el potencial como esta brecha es PlantVillage Nuru, una
aplicación móvil que ejecuta sin conexión un modelo de detección de síntomas foliares en
mandioca~\cite{ramcharan2019mobile}. En una evaluación de campo en Tanzania, Nuru alcanzó un
65\,\% de accuracy en el diagnóstico de enfermedades virales, por encima de los extensionistas
agrícolas (40--58\,\%) y de los productores (18--31\,\%), y llegó al 74--88\,\% cuando el
diagnóstico combinaba seis hojas por planta~\cite{mrisho2020accuracy}. Estos resultados
muestran que un modelo liviano ejecutado en el teléfono puede aportar valor real, pero
también que el desempeño en el campo es muy inferior al que se obtiene en el laboratorio.

En este trabajo desarrollamos un prototipo de clasificación de enfermedades foliares en
tres cultivos de importancia hortícola (pimiento, papa y tomate), con 15 clases que incluyen
hojas sanas y enfermas de PlantVillage. Partimos de modelos preentrenados en ImageNet y
priorizamos arquitecturas livianas (MobileNetV3~\cite{howard2019searching} y
EfficientNet-B0~\cite{tan2019efficientnet}), compatibles con una inferencia en CPU o en
dispositivos móviles. Las contribuciones principales son:
\begin{itemize}
  \item un protocolo de partición estratificada y agrupada por imágenes casi duplicadas, que
        reduce la fuga de datos entre entrenamiento y test;
  \item una comparación sistemática entre \emph{linear probe} y \emph{fine-tuning} parcial de
        distintos backbones, con y sin data augmentation;
  \item un análisis cuantitativo de explicabilidad con Grad-CAM~\cite{selvaraju2017gradcam},
        que mide qué fracción de la atención del modelo recae sobre la hoja y no sobre el fondo;
  \item una evaluación de robustez ante perturbaciones de la imagen y sobre datos externos a
        PlantVillage, junto con la medición del costo de inferencia en CPU.
\end{itemize}

El resto del trabajo se organiza de la siguiente manera: la Sección~\ref{sec:metodologia}
describe los datos, la partición, las arquitecturas y el protocolo experimental; la
Sección~\ref{sec:resultados} presenta y analiza los resultados; y la
Sección~\ref{sec:conclusiones} resume las conclusiones y las líneas de trabajo futuro.
